%TL;DR: Prior work studies activation control, multilingual safety, and agentic evaluation, but rarely combines role-specific intervention placement with prompt-matched multilingual outcome-composition analysis.
\section{Related Work}
\paragraph{Activation steering.} Activation-level interventions modify model behavior by adding directions derived from contrasts, probes, or behavioral examples. Prior work has studied refusal directions and the reliability of steering vectors across behaviors and prompt distributions. Our evaluation differs in three respects: it applies the intervention at distinct roles in a structured pipeline, preserves language-specific paired comparisons, and separates explicit decisions from unresolved outputs. We make no localization claim and do not treat behavioral contrasts as evidence of a specific circuit.

\paragraph{Multilingual safety.} Existing multilingual safety studies document uneven refusal and compliance behavior across languages and show that low-resource languages can expose alignment gaps. These studies motivate our language-stratified design. We extend descriptive baselining by evaluating prompt-matched intervention contrasts and by reporting whether effects survive predeclared correction families rather than extrapolating from one favourable language.

\paragraph{Multi-agent and agentic safety.} Agentic benchmarks evaluate tool use, harmful task completion, and guardrail behavior, often at the level of the complete system. Our pipeline-level analysis treats the intervention location as an experimental factor and asks whether an upstream Planner and downstream Critic respond differently to the same nominal steering operation.

\paragraph{Evaluation and statistical testing.} We use exact paired McNemar tests for binary matched outcomes, paired bootstrap percentile intervals for effect uncertainty, and Holm correction for predeclared families. This design follows standard paired evaluation practice while making correction families and decision-interface failures explicit.